\chapter{Efficienza Algoritmica sulle Architetture CPU Moderne}
\label{chap:efficienza-cpu}

\FloatBarrier
\section{Il Ruolo dell'Efficienza nell'Information Retrieval}
\label{sec:ruolo-efficienza-ir}

Lo sviluppo di sistemi di reperimento dell'informazione su scala globale si articola lungo due dimensioni complementari e in costante tensione tra loro~\cite{manning2008,buttcher2010}:

\begin{enumerate}
    \item \textbf{Efficacia (\textit{Effectiveness}):} massimizza la qualit\`a semantica dei risultati restituiti e la soddisfazione dell'utente (\textit{User Satisfaction}), adoperando modelli di ranking sofisticati, reti neurali e tecniche avanzate di Natural Language Processing tipicamente sviluppate in ambienti ad alto livello di astrazione (Python, PyTorch).
    \item \textbf{Efficienza (\textit{Efficiency}):} minimizza i costi computazionali, i consumi energetici e l'infrastruttura hardware necessaria per servire milioni di interrogazioni al secondo, ottimizzando il codice fino all'ultimo ciclo di clock della CPU e al singolo bit di memoria occupato.
\end{enumerate}

Chi progetta modelli ad alto livello non pu\`o ignorare i vincoli fisici dell'hardware sottostante per non creare soluzioni irrealizzabili in produzione; parallelamente, chi ottimizza a basso livello deve comprendere la semantica dei dati per preservare l'efficacia del recupero.

\FloatBarrier
\subsection{I Tre Pilastri dell'Efficienza nei Sistemi IR}
\label{subsec:tre-pilastri-efficienza}

L'efficienza algoritmica nei moderni sistemi di reperimento dell'informazione non \`e un mero esercizio di micro-ottimizzazione, ma un requisito abilitante che si fonda su tre pilastri:

\begin{enumerate}
    \item \textbf{Esperienza Utente (\textit{User Experience}):}
    Nei motori di ricerca Web contemporanei, la risposta a una query deve essere erogata entro una frazione minima di secondo (latenze tipiche inferiori a $50\text{--}100\,\text{ms}$). Ritardi percepibili inducono frustrazione e spingono l'utente a migrare verso servizi concorrenti. Inoltre, ridurre la latenza di computazione consente di eseguire algoritmi di ranking e filtraggio pi\`u complessi ed efficaci \emph{a parit\`a di budget temporale prestabilito} (\textit{Time Budget}).
    \item \textbf{Scalabilit\`a Operativa (\textit{Scalability}):}
    Le collezioni Web e aziendali contano decine di miliardi di documenti e gestiscono picchi di centinaia di migliaia di query al secondo. Il fabbisogno infrastrutturale di un motore di ricerca \`e approssimabile dalla relazione:
    \begin{equation}
    \text{Numero di Macchine} = \left\lceil \frac{\text{Volume Totale dei Dati e delle Query}}{\text{Capacit\`a Operativa del Singolo Nodo}} \right\rceil
    \end{equation}
    Un algoritmo scalabile richiede una crescita lineare dei nodi al raddoppiare dei dati; un algoritmo inefficiente rischia invece una crescita superlineare insostenibile. Poiché i dati operativi devono risiedere in RAM ed essere replicati su molteplici macchine per garantire la tolleranza ai guasti (\textit{fault tolerance}), qualsiasi inefficienza di memoria o di calcolo viene amplificata dal fattore di replica dell'infrastruttura.
    \item \textbf{Sostenibilit\`a Ambientale ed Economica (\textit{Environmental Sustainability}):}
    I moderni data center assorbono quantitativi colossali di energia elettrica, sia per alimentare i processori ad altissima densit\`a sia per i complessi impianti di dissipazione termica (\textit{cooling systems}). Un risparmio marginale del $5\%$ nell'utilizzo della CPU si traduce in milioni di kilowattora risparmiati all'anno e in un abbattimento diretto dell'impronta di carbonio (\textit{carbon footprint}).
\end{enumerate}

\FloatBarrier
\section{La Gerarchia di Memoria e i Principi di Localit\`a}
\label{sec:gerarchia-memoria-localita}

Nei sistemi di elaborazione contemporanei, le prestazioni reali sono dominate dai tempi di accesso alla memoria piuttosto che dalla mera velocit\`a aritmetica dell'ALU: si parla storicamente di \textbf{Memory Wall}~\cite{bakhvalov2024,slotin2023}.

\FloatBarrier
\subsection{Il Divario di Latenza tra CPU e DRAM}
\label{subsec:divario-latenza-cpu-dram}

Mentre il modello teorico astratto della macchina RAM assume un costo di accesso uniforme e costante $\BigO{1}$ per qualsiasi indirizzo, la fisica del silicio impone un divario abissale tra la velocit\`a di commutazione dei registri interni del processore e i tempi di risposta della memoria dinamica (DRAM). Per mitigare tale asimmetria, l'architettura dei calcolatori interpone una gerarchia di memorie cache via via pi\`u veloci, piccole e costose (Tabella~\ref{tab:gerarchia-memoria}).

\begin{table}[H]
\centering
\small
\resizebox{0.96\linewidth}{!}{%
\begin{tabular}{l r r r}
\toprule
\textbf{Livello Hardware} & \textbf{Latenza Tipica (ns)} & \textbf{Cicli CPU Equivalenti} & \textbf{Capacit\`a Tipica} \\
\midrule
Istruzione CPU (Registri / ALU) & $0.5\text{--}1.0\,\text{ns}$ & $\approx 1\text{--}2$ & $< 1\,\text{KB}$ (interni al core) \\
L1I Cache (Instruction) & $0.5\,\text{ns}$ & $\approx 1$ ciclo & $32\,\text{KB}$ per core \\
L1D Cache (Data) & $0.5\,\text{ns}$ & $\approx 1$ ciclo & $32\,\text{KB}$ per core \\
L2 Cache (Dati + Istruzioni) & $1.5\,\text{ns}$ & $\approx 3\text{--}4$ cicli & $256\,\text{KB}\text{--}1\,\text{MB}$ per core \\
L3 Cache (LLC, Condivisa) & $5.0\,\text{ns}$ & $\approx 12\text{--}15$ cicli & $16\text{--}64\,\text{MB}$ per package \\
RAM Principale (DRAM) & $80.0\,\text{ns}$ & $\approx 160\text{--}200+$ cicli & $64\text{--}512\,\text{GB}$ \\
Storage Secondario (SSD NVMe / Disco) & $500\,000\,\text{ns}$ ($0.5\,\text{ms}$) & oltre un milione & $8\text{--}30\,\text{TB}$ \\
\bottomrule
\end{tabular}%
}
\caption{Confronto quantitativo di latenze, cicli di clock e capacit\`a nella gerarchia di memoria di un moderno processore server x86\_64~\cite{bakhvalov2024,slotin2023}.}
\label{tab:gerarchia-memoria}
\end{table}

\begin{definition}[Algoritmo \textit{Memory-Bound} vs. \textit{CPU-Bound}]
Un algoritmo si definisce \textbf{memory-bound} quando il suo tempo di esecuzione totale \`e dominato dall'attesa del trasferimento dati dalla DRAM (stallo della CPU durante i cache miss), rendendo trascurabile il numero di operazioni matematiche eseguite. Si definisce \textbf{CPU-bound} quando tutti i dati necessari risiedono nelle cache ad alta velocit\`a (L1/L2) e il fattore limitante \`e la capacit\`a di calcolo o di throughput delle istruzioni del processore.
\end{definition}

\FloatBarrier
\subsection{Architettura Fisica sul Chip e Disaccoppiamento L1I / L1D}
\label{subsec:architettura-chip-l1}

Per evitare conflitti strutturali sul bus di memoria durante l'esecuzione parallela, la cache di primo livello \`e fisicamente sdoppiata:
\begin{itemize}
    \item \textbf{L1I (\textit{Instruction Cache}):} contiene unicamente il codice macchina prelevato dall'unit\`a di \textit{instruction fetch}.
    \item \textbf{L1D (\textit{Data Cache}):} serve le richieste di lettura e scrittura degli operandi (\texttt{load} e \texttt{store}) emesse dalle ALU.
\end{itemize}
La cache L2 \`e anch'essa integrata all'interno del singolo core (fornendo un buffer intermedio a bassa latenza), mentre la cache L3 (denominata \textit{Last Level Cache}, LLC) risiede sul die del chip ma \`e condivisa dinamicamente tra tutti i core fisici del socket.

\FloatBarrier
\subsection{La Cache Line come Unit\`a Atomica di Trasferimento}
\label{subsec:cache-line}

\begin{property}[Unit\`a di Trasferimento Hardware: $B = 64\text{ Byte}$]
I moderni processori non leggono né scrivono mai singoli byte o singole variabili isolate dalla memoria principale: l'unit\`a atomica di trasferimento tra DRAM e cache \`e la \textbf{cache line}, la cui dimensione universale nelle architetture x86\_64 ed ARM \`e pari a $B = 64\,\text{byte}$ (corrispondente a 8 interi a 64 bit).
\end{property}

Quando un'istruzione accede a un indirizzo di memoria $A$, l'hardware carica obbligatoriamente l'intero blocco allineato:
\begin{equation}
\left[\left\lfloor \frac{A}{64} \right\rfloor \times 64, \quad \left\lfloor \frac{A}{64} \right\rfloor \times 64 + 63\right]
\end{equation}
Se l'algoritmo legge un singolo intero a 64 bit ($8\,\text{byte}$) ed effettua un salto casuale verso un'altra area della memoria, i restanti $56\,\text{byte}$ (pari all'$87.5\%$ del traffico di bus assorbito) vengono irrimediabilmente sprecati.

\FloatBarrier
\subsection{Localit\`a Temporale, Spaziale e Hardware Prefetching}
\label{subsec:localita-prefetching}

Un algoritmo sfrutta la gerarchia di memoria se e solo se esibisce una delle due propriet\`a di localit\`a~\cite{bakhvalov2024}:
\begin{itemize}
    \item \textbf{Localit\`a Temporale:} Quando una locazione di memoria viene acceduta, \`e altamente probabile che la stessa locazione venga richiesta nuovamente a breve termine. Le cache sfruttano questo principio implementando politiche di rimpiazzo (varianti di \textbf{LRU}, \textit{Least Recently Used}) che mantengono in memoria i blocchi pi\`u recentemente referenziati.
    \item \textbf{Localit\`a Spaziale:} Quando una locazione viene acceduta, \`e altamente probabile che locazioni adiacenti vengano accedute nell'immediato futuro. Leggere sequenzialmente elementi contigui ammortizza il costo di caricamento della cache line su tutti i dati presenti nel blocco.
\end{itemize}

I processori integrano circuiti dedicati, detti \textbf{Hardware Prefetcher}~\cite{bakhvalov2024,slotin2023}, che analizzano a runtime il flusso degli indirizzi richiesti per anticipare i trasferimenti di memoria:
\begin{itemize}
    \item \textbf{Stream Prefetcher:} Rileva accessi puramente lineari ($A, A+64, A+128, \dots$) e ordina preventivamente alla DRAM il caricamento delle linee successive nelle cache L1/L2, mascherando completamente la latenza di caricamento.
    \item \textbf{Stride Prefetcher:} Intercetta accessi a passo costante ($A, A+s, A+2s, \dots$), caricando in anticipo le linee distanziate di $s$ byte.
    \item \textbf{Pattern Casuali (\textit{Pointer Chasing}):} Di fronte a salti casuali e non deterministici, nessun circuito predittivo \`e in grado di anticipare l'indirizzo bersaglio; la CPU \`e costretta ad arrestarsi a ogni iterazione, subendo per intero la latenza di $80\,\text{ns}$ della RAM.
\end{itemize}

\FloatBarrier
\section{Caso di Studio: Scansione Sequenziale contro \textit{Pointer Chasing}}
\label{sec:caso-studio-pointer-chasing}

Per quantificare sperimentalmente l'impatto della localit\`a e del prefetching, si esamina il seguente esperimento canonico. Si consideri un vettore \texttt{v} di dimensione $N$ contenente una \textbf{permutazione casuale con un singolo ciclo}:

\begin{definition}[Permutazione con un Singolo Ciclo]
Una permutazione $\pi$ dell'insieme $\{0, 1, \dots, N-1\}$ possiede un singolo ciclo se e solo se, partendo dall'indice iniziale $0$, la catena ricorsiva $k_{t+1} = \pi(k_t)$ visita ogni singola cella del vettore esattamente una volta prima di ritornare a $0$:
\begin{equation}
0 \to \pi(0) \to \pi(\pi(0)) \to \dots \to \pi^{N-1}(0) \to 0
\end{equation}
senza generare sottocicli disgiunti o componenti isolate.
\end{definition}

Si pongono a confronto due funzioni Rust che sommano tutti gli elementi di \texttt{v}:

\begin{lstlisting}[style=rustcode,caption={Scansione sequenziale con localit\`a spaziale e prefetching.}]
pub fn sum(v: &[usize]) -> usize {
    let mut sum = 0;
    for &value in v.iter() {
        sum += value;
    }
    sum
}
\end{lstlisting}

\begin{lstlisting}[style=rustcode,caption={Scansione disordinata per salti casuali (\textit{pointer chasing}).}]
pub fn sum_by_jumping(v: &[usize]) -> usize {
    let mut sum = 0;
    let mut k = 0;
    for _ in 0..v.len() {
        sum += v[k];
        k = v[k]; // il valore letto determina l'indirizzo del passo successivo!
    }
    sum
}
\end{lstlisting}

Nel modello computazionale teorico RAM entrambe le routine eseguono esattamente $N$ addizioni e $N$ accessi a vettore, presentando identica complessit\`a asintotica $\BigO{N}$. Sulla CPU reale, misurando il tempo medio per singolo elemento al crescere della taglia $N$, emergono divergenze impressionanti (Tabella~\ref{tab:sum-vs-jumping}).

\begin{table}[H]
\centering
\small
\resizebox{0.96\linewidth}{!}{%
\begin{tabular}{r r l r r r}
\toprule
\textbf{Elementi ($N$)} & \textbf{Dimensione} & \textbf{Livello Cache Ospitante} & \textbf{\texttt{sum()}} & \textbf{\texttt{sum\_by\_jumping()}} & \textbf{Rallentamento} \\
\midrule
$256$ & $2\,\text{KiB}$ & Cache L1 ($32\,\text{KB}$) & $0.1\,\text{ns}$ & $1.3\,\text{ns}$ & \textbf{$13\times$} \\
$8\,192$ & $64\,\text{KiB}$ & Cache L2 ($256\,\text{KB}$) & $0.1\,\text{ns}$ & $2.5\,\text{ns}$ & \textbf{$25\times$} \\
$262\,144$ & $2\,\text{MiB}$ & Cache L3 ($16\,\text{MB}$) & $0.2\,\text{ns}$ & $13.4\,\text{ns}$ & \textbf{$67\times$} \\
$16\,777\,216$ & $128\,\text{MiB}$ & Memoria RAM (DRAM) & $0.4\,\text{ns}$ & $84.9\,\text{ns}$ & \textbf{$212\times$} \\
\bottomrule
\end{tabular}%
}
\caption{Benchmark comparativo del tempo per elemento al variare della taglia del vettore~\cite{bakhvalov2024}.}
\label{tab:sum-vs-jumping}
\end{table}

\FloatBarrier
\subsection{Strategie di Ottimizzazione nel Pointer Chasing}
\label{subsec:sfide-didattiche-pointer-chasing}

\begin{example}[Challenge 2: Sfruttare la Cache Line nel \textit{Pointer Chasing}]
Come ammortizzare l'attesa di $85\,\text{ns}$ della RAM effettuando comunque salti disordinati?
Se raggruppiamo gli elementi in blocchi da 64 byte (8 interi \texttt{usize}), possiamo saltare casualmente da un blocco all'altro, ma accumulare tutti gli 8 valori della linea prima di passare al blocco successivo:
\begin{lstlisting}[style=rustcode]
pub fn sum_by_jumping_cache_lines(v: &[usize]) -> usize {
    let mut sum = 0;
    let mut k = 0;
    let num_blocks = v.len() / 8;
    for _ in 0..num_blocks {
        let base = k * 8;
        // La prima lettura carica forzatamente tutti i 64 byte in L1
        sum += v[base] + v[base+1] + v[base+2] + v[base+3]
             + v[base+4] + v[base+5] + v[base+6];
        k = v[base+7] % num_blocks; // puntatore al prossimo blocco casuale
    }
    sum
}
\end{lstlisting}
Il cache miss di $85\,\text{ns}$ viene pagato una sola volta per blocco e suddiviso su 8 addizioni: il costo scende a $\approx \frac{85\,\text{ns}}{8} \approx 10.6\,\text{ns}$ per elemento, con un'accelerazione di circa $8\times$ rispetto al salto singolo.
\end{example}

\begin{example}[Challenge 1: Eliminare il Lieve Rallentamento di \texttt{sum()} su Grandi Dati]
In Tabella~\ref{tab:sum-vs-jumping}, \texttt{sum()} passa da $0.1\,\text{ns}$ a $0.4\,\text{ns}$ quando $N=16\text{M}$ perch\'e l'hardware prefetcher, pur saturando la banda del bus, non riesce ad alimentare l'ALU a pieno regime (il tempo di calcolo di una somma \`e inferiore al tempo di trasferimento dalla DRAM).
Se arricchiamo il corpo del ciclo con operazioni aritmetiche pi\`u dense su ciascun elemento (ad esempio molteplici moltiplicazioni, polinomi o hashing), il tempo computazionale supera la latenza di caricamento ($T_{\text{comp}} \gg T_{\text{load}}$): il sistema diviene interamente \textbf{CPU-bound} e la curva dei tempi di \texttt{sum()} rimane perfettamente costante al variare di $N$.
\end{example}

\FloatBarrier
\section{Parallel Computing e Forme di Parallelismo nella CPU}
\label{sec:parallel-computing-cpu}

Per incrementare le prestazioni computazionali, l'architettura dei calcolatori moderni non si limita a innalzare la frequenza di clock, ma sfrutta molteplici forme di \textbf{Parallel Computing}~\cite{bakhvalov2024,slotin2023}.

\begin{noteblock}{Parallelismo a Livello di Server: Il Paradigma \textit{Embarrassingly Parallel}}
Nei motori di ricerca su larga scala, la gestione simultanea di migliaia di query non richiede complesse sincronizzazioni tra i core per una singola interrogazione: si adotta il \textbf{parallelismo triviale} (\textit{embarrassingly parallel}), assegnando ciascuna query a un thread isolato su un core dedicato. Non essendovi condivisione di memoria o contese di lock, si azzerano gli overhead di coordinamento e i problemi di coerenza della cache.
\end{noteblock}

All'interno del \textbf{singolo core}, la CPU sfrutta simultaneamente tre livelli di parallelismo microscopico:
\begin{enumerate}
    \item \textbf{Pipelining:} Sovrappone l'esecuzione temporale di istruzioni contigue, suddividendo il ciclo di vita dell'istruzione in stadi indipendenti.
    \item \textbf{Processori Superscalari:} Integrano molteplici unit\`a di calcolo parallele (\textit{Execution Ports}) all'interno del medesimo core per elaborare pi\`u istruzioni indipendenti nello stesso ciclo di clock ($IPC > 1$).
    \item \textbf{SIMD (\textit{Single Instruction, Multiple Data}):} Esegue la medesima operazione su un intero vettore di dati paralleli mediante registri dedicati di ampiezza fino a 512 bit.
\end{enumerate}

\FloatBarrier
\subsection{Anatomia della Pipeline a Cinque Stadi}
\label{subsec:pipeline-5-stadi}

Il pipelining scompone l'esecuzione di ogni istruzione macchina in una catena di montaggio sequenziale. Nel modello didattico classico si distinguono 5 stadi (Figura~\ref{fig:pipeline-5-stadi}):
\begin{enumerate}
    \item \textbf{IF (\textit{Instruction Fetch}):} Prelievo dell'istruzione binaria dalla cache L1I tramite il registro \textit{Program Counter} (PC).
    \item \textbf{ID (\textit{Instruction Decode}):} Decodifica dell'istruzione in $\mu$-ops e lettura degli operandi dai registri architetturali.
    \item \textbf{EX (\textit{Execute}):} Esecuzione del calcolo nelle unit\`a aritmetico-logiche (ALU) o calcolo dell'indirizzo effettivo di memoria.
    \item \textbf{MEM (\textit{Memory Access}):} Accesso alla cache L1D per istruzioni di caricamento (\texttt{load}) o scrittura (\texttt{store}).
    \item \textbf{WB (\textit{Write-Back}):} Scrittura del risultato definitivo nei registri di destinazione e ritiro (\textit{retirement}) dell'istruzione.
\end{enumerate}

\begin{figure}[H]
\centering
\begin{tikzpicture}[
    stage/.style={rectangle, draw=navy, fill=codebackground, thick, minimum width=1.35cm, minimum height=0.75cm, align=center, font=\scriptsize\bfseries},
    lbl/.style={font=\scriptsize, text=navy!80!black}
]
    \foreach \i/\name/\y in {0/i/0, 1/i+1/-0.9, 2/i+2/-1.8} {
        \node[lbl, anchor=east] at (-0.3, \y) {Istr. \name};
    }
    \foreach \s [count=\c from 0] in {IF,ID,EX,MEM,WB} {
        \node[stage] at (\c*1.5, 0) {\s};
    }
    \foreach \s [count=\c from 1] in {IF,ID,EX,MEM,WB} {
        \node[stage] at (\c*1.5, -0.9) {\s};
    }
    \foreach \s [count=\c from 2] in {IF,ID,EX,MEM,WB} {
        \node[stage] at (\c*1.5, -1.8) {\s};
    }
    \node[lbl] at (0, 0.65) {$t_1$};
    \foreach \t [count=\c from 2] in {2,...,7} {
        \node[lbl] at (\c*1.5-1.5, 0.65) {$t_{\c}$};
    }
\end{tikzpicture}
\caption{Esecuzione sovrapposta nella pipeline a 5 stadi: a ogni ciclo di clock $t$, ciascuno stadio hardware elabora una fase di un'istruzione differente.}
\label{fig:pipeline-5-stadi}
\end{figure}

\FloatBarrier
\subsection{Latenza, Throughput e Profondit\`a della Pipeline}
\label{subsec:latenza-throughput-pipeline}

Sia $k$ il numero di stadi della pipeline e $\tau$ il periodo del ciclo di clock.
\begin{itemize}
    \item La \textbf{Latenza} $L = k \cdot \tau$ \`e il tempo impiegato da una singola istruzione per attraversare l'intera catena di elaborazione.
    \item Il \textbf{Throughput} a regime $W = 1/\tau$ misura il numero di istruzioni ritirate per unit\`a di tempo: idealmente, a ogni ciclo di clock un'istruzione completa lo stadio di WB.
    \item Il periodo minimo di clock \`e imposto dal ritardo dello stadio circuitale pi\`u lento:
    \begin{equation}
    \tau \ge \max_{1 \le j \le k} \big(t_{\text{stadio}_j}\big)
    \end{equation}
\end{itemize}

Per aumentare la frequenza di clock (ridurre $\tau$ ed elevare i gigahertz), i progettisti aumentano la profondit\`a $k$ della pipeline (portandola tipicamente a $15\text{--}20$ stadi nei processori x86\_64 moderni come Intel Skylake o AMD Zen), suddividendo ogni operazione in micro-fasi circuitali estremamente brevi. Il prezzo da pagare per pipeline molto profonde risiede nella grave dilatazione della penalit\`a di svuotamento (\textit{pipeline flush}) in caso di errore di predizione dei salti.

\FloatBarrier
\section{Esecuzione Superscalare e Interleaving Aritmetico}
\label{sec:esecuzione-superscalare}

Un processore si definisce \textbf{superscalare} quando ha la capacit\`a di emettere, eseguire e ritirare pi\`u di un'istruzione per ciclo di clock all'interno del medesimo core ($IPC > 1$, \textit{Instructions Per Cycle})~\cite{bakhvalov2024}. 

I core ad alte prestazioni integrano multiple unit\`a funzionali parallele collegate a una serie di porte di esecuzione indipendenti (\textit{Execution Ports}):
\begin{itemize}
    \item Multiple ALU aritmetico-logiche per operazioni su interi e logica binaria;
    \item Moltiplicatori e divisori hardware dedicati;
    \item Porte di generazione degli indirizzi e unit\`a di caricamento/scrittura (\texttt{load}/\texttt{store});
    \item Unit\`a vettoriali SIMD per calcolo in virgola mobile e matriciale.
\end{itemize}

\FloatBarrier
\subsection{L'Esperimento dell'\textit{Interleaving} Aritmetico}
\label{subsec:esperimento-interleaving}

Si consideri un set di benchmark sperimentale in cui si accumulano riduzioni su un vettore di interi a 64 bit \texttt{\&[u64]}:

\begin{lstlisting}[style=rustcode,caption={Funzioni di benchmark: incremento delle operazioni indipendenti nel ciclo.}]
pub fn do_stuff(v: &[u64]) -> u64 {
    let mut prod = 1;
    for &value in v.iter() {
        prod *= value;
    }
    prod
}

pub fn do_more_stuff_2(v: &[u64]) -> (u64, u64, u64) {
    let mut prod = 1;
    let mut sum = 0;
    let mut xor = 0;
    for &value in v.iter() {
        prod *= value;
        sum += value;
        xor ^= value;
    }
    (prod, sum, xor)
}

#[inline]
pub fn do_more_stuff_3(v: &[u64]) -> (u64, u64, u64, u64) {
    let mut prod = 1;
    let mut sum = 0;
    let mut xor = 0;
    let mut sum2 = 0;
    for &value in v.iter() {
        prod *= value;
        sum += value;
        xor ^= value;
        sum2 += value;
    }
    (prod, sum, xor, sum2)
}
\end{lstlisting}

I risultati sperimentali mostrano un fenomeno controintuitivo (Tabella~\ref{tab:interleaving-superscalar}):

\begin{table}[H]
\centering
\small
\resizebox{\columnwidth}{!}{%
\begin{tabular}{l r r}
\toprule
\textbf{Funzione} & \textbf{Operazioni Eseguite per Iterazione} & \textbf{Tempo Medio per Iterazione} \\
\midrule
\texttt{do\_stuff()} & 1 moltiplicazione & $0.4\text{--}1.1\,\text{ns}$ \\
\texttt{do\_more\_stuff\_2()} & 3 operazioni (moltiplicazione, somma, XOR) & $\mathbf{0.4\,\text{ns}}$ \\
\texttt{do\_more\_stuff\_3()} & 4 operazioni (moltiplicazione, 2 somme, XOR) & $\mathbf{0.5\,\text{ns}}$ \\
\bottomrule
\end{tabular}%
}
\caption{Tempi per iterazione: eseguire tre o quattro calcoli indipendenti richiede lo stesso tempo di uno solo grazie all'architettura superscalare.}
\label{tab:interleaving-superscalar}
\end{table}

\textbf{Spiegazione fisica:}
Le variabili di accumulo \texttt{prod}, \texttt{sum} e \texttt{xor} sono completamente disaccoppiate tra loro. La CPU superscalare rileva l'assenza di dipendenze e invia contemporaneamente la moltiplicazione a una porta di calcolo complessa, la somma a una prima porta ALU e lo XOR a una seconda porta ALU. Le tre operazioni vengono completate \textbf{in parallelo nello stesso identico ciclo di clock}: il lavoro computazionale triplica, ma il tempo di esecuzione resta invariato ($0.4\,\text{ns}$).

\FloatBarrier
\section{Pipeline Hazards: Tipologie e Risoluzione}
\label{sec:pipeline-hazards}

Un \textbf{Pipeline Hazard} \`e una condizione architetturale che impedisce all'istruzione successiva di iniziare la propria esecuzione nel ciclo di clock immediatamente seguente, costringendo il processore a inserire bolle di stallo (\textit{pipeline bubbles} o \textit{stalls}) che degradano il throughput~\cite{bakhvalov2024,slotin2023}.

Si distinguono tre categorie fondamentali di hazard:

\FloatBarrier
\subsection{Hazard Strutturali (\textit{Structural Hazards})}
\label{subsec:structural-hazards}
Si verificano quando due o pi\`u istruzioni pronte per l'esecuzione richiedono simultaneamente la medesima risorsa hardware fisica (ad esempio una specifica porta ALU o un'unit\`a di divisione in virgola mobile non completamente pipelined). Rappresentano limiti fisici invalicabili del silicio, non risolvibili se non dilatando l'hardware.

\FloatBarrier
\subsection{Data Hazards (\textit{Dipendenze di Dato, RAW})}
\label{subsec:data-hazards}
Si verificano quando un'istruzione necessita del valore prodotto da un'istruzione precedente che non ha ancora completato il proprio stadio di esecuzione o di scrittura nei registri (\textit{Read-After-Write}, RAW). 

Si analizzi la funzione \texttt{do\_more\_stuff\_4()}, che introduce una rigida catena seriale di dipendenze:

\begin{lstlisting}[style=rustcode,caption={Catena di dipendenze seriali: ogni operazione dipende dal risultato della precedente.}]
pub fn do_more_stuff_4(v: &[u64]) -> (u64, u64, u64) {
    let mut prod = 1;
    let mut sum = 0;
    let mut xor = 0;
    for &value in v.iter() {
        sum += value;       // Calcola 'sum'
        prod *= sum;        // ATTENDE 'sum'!
        xor ^= prod;        // ATTENDE 'prod'!
        sum += xor;         // ATTENDE 'xor'!
    }
    (prod, sum, xor)
}
\end{lstlisting}

\begin{alertblock}{Il Crollo di Throughput causato dai Data Hazard}
Nella funzione \texttt{do\_more\_stuff\_4()}, la moltiplicazione non pu\`o iniziare finché la somma non \`e pronta; lo XOR non pu\`o iniziare prima del completamento della moltiplicazione; e la somma finale deve attendere lo XOR. 
La pipeline non pu\`o sovrapporre i calcoli: il tempo per iterazione balza da $\mathbf{0.4\,\text{ns}}$ a $\mathbf{1.7\,\text{ns}}$ (un rallentamento di oltre $4\times$).

\textbf{Soluzione Algoritmica:} Ristrutturare l'algoritmo spezzando le catene critiche di dipendenza o applicando l'\textit{interleaving} con operazioni disaccoppiate.
\end{alertblock}

\FloatBarrier
\subsection{Control Hazards (\textit{Dipendenze di Controllo})}
\label{subsec:control-hazards}
Si verificano in corrispondenza di istruzioni di salto condizionato (\texttt{if}, salti assembly \texttt{JNE}, \texttt{JLE}), in cui il processore non pu\`o determinare quale blocco di codice caricare finché la condizione logica non viene calcolata nello stadio EX. Di fronte a questo dilemma, la CPU \`e costretta a speculare.

\FloatBarrier
\section{Branch Prediction ed Esecuzione Speculativa}
\label{sec:branch-prediction}

Per evitare che la pipeline si arresti a ogni salto condizionato, i processori integrano un componente ad altissima complessit\`a: il \textbf{Branch Predictor}~\cite{bakhvalov2024,slotin2023}.

Il predittore prova a indovinare quale ramo verr\`a intrapreso prima ancora che la condizione sia valutata:
\begin{itemize}
    \item \textbf{Esecuzione Speculativa:} La CPU preleva ed esegue provvisoriamente le istruzioni del ramo predetto.
    \item \textbf{Predizione Corretta:} L'esecuzione procede a piena velocit\`a senza alcuno stallo.
    \item \textbf{Errore di Predizione (\textit{Branch Misprediction}):} Se la predizione si rivela errata, la CPU deve cancellare tutti i risultati provvisori calcolati ed eseguire il \textbf{flush completo della pipeline}, ricaricando le istruzioni corrette dal punto di salto.
\end{itemize}

\begin{property}[Costo del Branch Misprediction]
La penalit\`a in cicli di clock di un errore di predizione \`e pari al numero di stadi della pipeline tra il fetch (IF) e la verifica dell'istruzione nello stadio di esecuzione (EX), tipicamente \textbf{15--20 cicli di clock sprecati} per ogni singolo errore!
\end{property}

\FloatBarrier
\subsection{L'Esperimento dei Dati Ordinati e Alternati}
\label{subsec:esperimento-branch-prediction}

Si analizzi la funzione \texttt{do\_more\_stuff\_5()}, che inserisce un salto condizionato all'interno del ciclo:

\begin{lstlisting}[style=rustcode,caption={Presenza di salto condizionato all'interno del loop.}]
pub fn do_more_stuff_5(v: &[u64]) -> (u64, u64, u64) {
    let mut prod = 1;
    let mut sum = 0;
    let mut xor = 0;
    for &value in v.iter() {
        if value > 50 { // Salto condizionale!
            sum += value;
            prod *= sum;
            xor ^= prod;
            sum += xor;
        }
    }
    (prod, sum, xor)
}
\end{lstlisting}

Eseguendo la funzione su un vettore contenente valori estratti uniformemente in $[0, 99]$ sotto tre diverse disposizioni, emergono i risultati di Tabella~\ref{tab:branch-prediction-benchmark}.

\begin{table}[H]
\centering
\small
\resizebox{\columnwidth}{!}{%
\begin{tabular}{l r p{8.5cm}}
\toprule
\textbf{Disposizione dei Dati in \texttt{v}} & \textbf{Tempo per Iterazione} & \textbf{Comportamento del Branch Predictor} \\
\midrule
\texttt{do\_more\_stuff\_4()} (nessun \texttt{if}) & $1.7\,\text{ns}$ & Assenza di salti condizionati \\
\texttt{do\_more\_stuff\_5()} (dati casuali) & $\mathbf{3.5\,\text{ns}}$ & Ramo impredicibile: fallimento nel $\approx 50\%$ dei casi con continui flush \\
\texttt{do\_more\_stuff\_5()} (dati ordinati) & $\mathbf{1.2\,\text{ns}}$ & Predizione perfetta: solo 1 misprediction al punto di transizione $> 50$ \\
\texttt{do\_more\_stuff\_5()} (dati alternati) & $\mathbf{1.1\,\text{ns}}$ & Pattern periodico riconosciuto dai registri di shift storico: 0\% misprediction \\
\bottomrule
\end{tabular}%
}
\caption{Tempo medio per iterazione di \texttt{do\_more\_stuff\_5()} in funzione della predicibilit\`a dei dati~\cite{slotin2023}.}
\label{tab:branch-prediction-benchmark}
\end{table}

\begin{itemize}
    \item Con \textbf{dati casuali}, il salto \`e analogo al lancio di una moneta: la CPU sbaglia la met\`a delle volte, pagando costantemente la penalit\`a di $15\text{--}20$ cicli di flush ($\implies 3.5\,\text{ns}$).
    \item Con \textbf{dati ordinati}, il predittore assume che il salto non sia preso per tutta la prima met\`a e sia sempre preso per la seconda: commette un solo errore nell'intero vettore, impiegando solo $1.2\,\text{ns}$ (meno della versione senza condizione, poiché met\`a degli elementi viene saltata a costo zero!).
    \item Con \textbf{dati alternati} (un valore $>50$, uno $\le 50$), i predittori hardware integrano tabelle storiche dei pattern a due livelli (\textit{Branch History Shift Registers}): identificata l'alternanza periodica, il circuito anticipa la scelta con precisione assoluta ($\implies 1.1\,\text{ns}$).
\end{itemize}

\FloatBarrier
\subsection{Tecniche \textit{Branchless} e Condizione del Ciclo \texttt{for}}
\label{subsec:branchless-for}

Quando i dati non sono ordinabili né prevedibili, l'unica soluzione per evitare i disastrosi stalli di misprediction consiste nell'eliminare fisicamente il salto condizionato (\textbf{Branchless Programming}), trasformando la selezione logica in un'espressione aritmetica:
\begin{lstlisting}[style=rustcode]
// Versione branchless: converte il predicato booleano in 1 oppure 0
let mask = (value > 50) as u64;
sum += value * mask;
\end{lstlisting}
Il compilatore traduce l'espressione in istruzioni condizionali prive di salto (quali \texttt{CMOV} o istruzioni vettoriali mascherate): il processore esegue sempre il calcolo, ma la pipeline procede in modo lineare e privo di flush.

\begin{remark}[Perch\'e il Salto di Fine Ciclo \texttt{for} non Danneggia le Prestazioni]
Ogni ciclo \texttt{for} termina con un'istruzione di salto condizionato per ritornare all'inizio. Questo salto non rallenta la CPU perch\'e \`e intrinsecamente sbilanciato: su $N$ iterazioni il salto viene preso $N-1$ volte consecutive ed \`e non preso solo all'uscita finale. Il branch predictor apprende istantaneamente la regolarit\`a, commettendo un unico errore su milioni di passi.
\end{remark}

\FloatBarrier
\section{Vettorizzazione SIMD e la Risoluzione del Mistero di L1 Cache}
\label{sec:simd-vettorizzazione}

La terza forma cardine di parallelismo interno alla CPU \`e rappresentata dal paradigma \textbf{SIMD} (\textit{Single Instruction, Multiple Data})~\cite{bakhvalov2024,slotin2023}.

\FloatBarrier
\subsection{Registri Vettoriali e Famiglie di Estensioni x86}
\label{subsec:registri-vettoriali}

Le moderne CPU supportano estensioni del set di istruzioni progettate per operare su registri speciali ad altissima capacit\`a:
\begin{itemize}
    \item \textbf{MMX / SSE (Streaming SIMD Extensions):} Registri \texttt{XMM} a 128 bit.
    \item \textbf{AVX / AVX2 (Advanced Vector Extensions):} Registri \texttt{YMM} a 256 bit.
    \item \textbf{AVX-512:} Registri \texttt{ZMM} a 512 bit (in grado di alloggiare contemporaneamente 64 interi a 8 bit o 8 numeri a virgola mobile a 64 bit).
\end{itemize}

I registri vettoriali vengono idealmente suddivisi in corsie parallele (\textit{vector lanes}): una singola istruzione macchina (ad esempio \texttt{vpaddq}) comanda l'esecuzione simultanea della medesima operazione algebrica su tutti gli elementi delle corsie parallele, moltiplicando linearmente il throughput computazionale (Figura~\ref{fig:simd-concept}).

\begin{figure}[H]
\centering
\resizebox{0.95\columnwidth}{!}{%
\begin{tikzpicture}[
    box/.style={rectangle, draw=navy, fill=codebackground, thick, minimum width=2.1cm, minimum height=0.85cm, font=\footnotesize\bfseries, align=center, rounded corners=3pt},
    op/.style={circle, draw=darkgreen, fill=darkgreen!15, thick, inner sep=3pt, font=\small\bfseries}
]
    % Titoli delle 4 corsie (Lanes)
    \node[font=\footnotesize\bfseries, text=navy] at (0.8, 3.4) {Corsia 0 [63:0]};
    \node[font=\footnotesize\bfseries, text=navy] at (3.6, 3.4) {Corsia 1 [127:64]};
    \node[font=\footnotesize\bfseries, text=navy] at (6.4, 3.4) {Corsia 2 [191:128]};
    \node[font=\footnotesize\bfseries, text=navy] at (9.2, 3.4) {Corsia 3 [255:192]};

    % Etichette di riga a sinistra con ampio distanziamento
    \node[anchor=east, font=\small\bfseries, align=right] at (-0.9, 2.5) {Registro A (\texttt{ymm0}):\\[2pt]\scriptsize 256 bit (4 $\times$ 64 bit)};
    \node[anchor=east, font=\small\bfseries, align=right] at (-0.9, 1.25) {Registro B (\texttt{ymm1}):\\[2pt]\scriptsize 256 bit (4 $\times$ 64 bit)};
    \node[anchor=east, font=\small\bfseries, align=right, text=darkgreen] at (-0.9, 0.0) {Operazione SIMD:\\[2pt]\scriptsize 4 ALU parallele (\texttt{vpaddq})};
    \node[anchor=east, font=\small\bfseries, align=right] at (-0.9, -1.35) {Registro C (\texttt{ymm2}):\\[2pt]\scriptsize Destinazione (256 bit)};

    % Registro A
    \node[box] (a0) at (0.8, 2.5) {$A_0$};
    \node[box] (a1) at (3.6, 2.5) {$A_1$};
    \node[box] (a2) at (6.4, 2.5) {$A_2$};
    \node[box] (a3) at (9.2, 2.5) {$A_3$};

    % Registro B
    \node[box] (b0) at (0.8, 1.25) {$B_0$};
    \node[box] (b1) at (3.6, 1.25) {$B_1$};
    \node[box] (b2) at (6.4, 1.25) {$B_2$};
    \node[box] (b3) at (9.2, 1.25) {$B_3$};

    % Operatori ALU (Somma SIMD)
    \node[op] (plus0) at (0.8, 0.0) {$+$};
    \node[op] (plus1) at (3.6, 0.0) {$+$};
    \node[op] (plus2) at (6.4, 0.0) {$+$};
    \node[op] (plus3) at (9.2, 0.0) {$+$};

    % Registro C (Destinazione)
    \node[box, draw=darkgreen, fill=darkgreen!12] (c0) at (0.8, -1.35) {$C_0 = A_0 + B_0$};
    \node[box, draw=darkgreen, fill=darkgreen!12] (c1) at (3.6, -1.35) {$C_1 = A_1 + B_1$};
    \node[box, draw=darkgreen, fill=darkgreen!12] (c2) at (6.4, -1.35) {$C_2 = A_2 + B_2$};
    \node[box, draw=darkgreen, fill=darkgreen!12] (c3) at (9.2, -1.35) {$C_3 = A_3 + B_3$};

    % Flussi per ciascuna corsia: A entra da sinistra, B entra dall'alto
    \foreach \i in {0,1,2,3} {
        \draw[-Stealth, thick, navy] (a\i.west) -- ++(-0.35, 0) |- (plus\i.west);
        \draw[-Stealth, thick, navy] (b\i.south) -- (plus\i.north);
        \draw[-Stealth, thick, darkgreen] (plus\i.south) -- (c\i.north);
    }
\end{tikzpicture}%
}
\caption{Schema concettuale di operazione SIMD vettoriale a 256 bit: quattro addizioni a 64 bit eseguite simultaneamente in un unico ciclo di clock.}
\label{fig:simd-concept}
\end{figure}

\FloatBarrier
\subsection{Auto-Vettorizzazione del Compilatore contro Intrinsics}
\label{subsec:autovettorizzazione-intrinsics}

Lo sfruttamento di SIMD nel software segue due strade:

\begin{enumerate}
    \item \textbf{Auto-vettorizzazione da parte del Compilatore:}
    Compilatori ottimizzanti moderni (LLVM/Rust, GCC, Clang con flag \texttt{-O3}) analizzano autonomamente i cicli for privi di effetti collaterali e traducono le sequenze scalari in istruzioni vettoriali SIMD.
    \begin{itemize}
        \item \textit{Vantaggi:} Totale trasparenza, nessun intervento manuale sul codice sorgente, portabilit\`a cross-platform immediata.
        \item \textit{Limiti:} Il compilatore \`e conservativo; se il ciclo contiene condizioni complesse, accessi indiretti o rischio di aliasing di puntatori, l'auto-vettorizzazione fallisce.
    \end{itemize}
    \item \textbf{Intrinsics Hardware di Basso Livello:}
    Funzioni speciali C/Rust (ad es. libreria \texttt{core::arch::x86\_64}) che mappano biunivocamente su specifiche istruzioni assembly SIMD (quali \texttt{\_mm256\_loadu\_si256}, \texttt{\_mm256\_add\_epi64}).
    \begin{itemize}
        \item \textit{Vantaggi:} Controllo totale sulle strutture dati e possibilit\`a di ottimizzare algoritmi complessi non identificati dal compilatore.
        \item \textit{Limiti:} Richiede una profonda padronanza micro-architetturale e rende il codice rigidamente vincolato allo specifico set di istruzioni della CPU bersaglio.
    \end{itemize}
\end{enumerate}

\FloatBarrier
\subsection{La Rivelazione Finale: Perch\'e \texttt{sum()} valeva $0.1\,\text{ns}$ in L1 Cache?}
\label{subsec:soluzione-enigma-l1}

A conclusione della trattazione sull'efficienza della CPU, si ricollega l'intero percorso didattico all'esperimento iniziale.

\begin{alertblock}{La Risoluzione dell'Enigma: L'Analisi con Compiler Explorer}
Nella \secref{sec:caso-studio-pointer-chasing}, per $N=256$ elementi (vettore di appena $2\,\text{KiB}$, interamente contenuto nella velocissima cache L1D):
\begin{center}
\texttt{sum()} impiegava $\mathbf{0.1\,\text{ns}}$ \quad contro i $\mathbf{1.3\,\text{ns}}$ di \texttt{sum\_by\_jumping()} ($13\times$ di differenza!)
\end{center}
Poiché entrambi i vettori risiedevano interamente in cache L1, il divario \textbf{non poteva dipendere dai cache miss della RAM}! Qual era dunque la causa fisica di tale differenza abissale?

Ispezionando l'assembly generato tramite \textbf{Compiler Explorer} (Godbolt), emerge la verit\`a:
\begin{enumerate}
    \item \textbf{\texttt{sum()} \`e stata auto-vettorizzata con AVX dal compilatore!}
    Il ciclo non compie addizioni scalari elemento per elemento: carica vettori a 256 bit e calcola \textbf{quattro somme da 64 bit simultaneamente in un singolo ciclo di clock} (\texttt{vpaddq}), srotola il loop (\textit{loop unrolling}) ed esegue una sola somma orizzontale di riduzione al termine.
    \item \textbf{\texttt{sum\_by\_jumping()} \`e strutturalmente non vettorizzabile!}
    La dipendenza di dato stretta $k_{t+1} = v[k_t]$ impone che l'indirizzo della lettura successiva dipenda dal valore caricato dalla memoria nell'iterazione corrente. Il compilatore \`e impossibilitato a vettorizzare: la CPU esegue un'unica somma scalare seriale per volta, subendo inoltre la latenza di dipendenza RAW.
\end{enumerate}
\end{alertblock}

Questo risultato sancisce il principio metodologico fondamentale di questo capitolo:
\begin{property}[Principio di Regime Operativo]
\begin{itemize}
    \item Quando i dati eccedono le cache e risiedono nella memoria principale DRAM, il programma \`e \textbf{memory-bound}: le prestazioni sono dominate dalla localit\`a spaziale, dalla dimensione della cache line ($64\,\text{byte}$) e dall'efficacia dell'hardware prefetcher.
    \item Quando i dati sono compressi o ridotti al punto da risiedere interamente nelle cache L1/L2, il programma diventa \textbf{CPU-bound}: le prestazioni sono dominate dall'assenza di stalli della pipeline, dalla predicibilit\`a dei salti e dalla capacit\`a di sfruttare la vettorizzazione SIMD.
\end{itemize}
\end{property}

\FloatBarrier
\section{Quadro Sinottico e Implicazioni per l'Information Retrieval}
\label{sec:quadro-sinottico-ir-hardware}

La Tabella~\ref{tab:sinottico-hardware-ir} riassume schematicamente tutti i principi hardware esaminati e illustra le dirette implicazioni per la progettazione degli algoritmi e delle strutture dati nei motori di ricerca moderni (temi che verranno sviluppati nei successivi Capitoli 6, 7 e 8).

\begin{table}[H]
\centering
\small
\resizebox{\columnwidth}{!}{%
\begin{tabular}{l l l}
\toprule
\textbf{Principio Hardware} & \textbf{Meccanismo / Vincolo Fisico} & \textbf{Impatto nella Progettazione di Sistemi IR} \\
\midrule
\textbf{Gerarchia di Memoria} & Divario RAM vs CPU ($80\,\text{ns}$ vs $0.5\,\text{ns}$) & Compressione degli indici invertiti per farli risiedere in RAM \\
\textbf{Cache Line ($64\text{ B}$)} & Blocco atomico minimo di trasferimento & Layout a blocchi contigui (\textit{Array of Structures} vs \textit{SoA}) \\
\textbf{Hardware Prefetcher} & Anticipazione di stream contigui o a passo fisso & Scansione sequenziale delle posting list anziché alberi sparsi \\
\textbf{Pipelining} & Sovrapposizione di istruzioni a stadi sequenziali & Funzioni interne snelle senza chiamate a funzione indirette \\
\textbf{Superscalarità} & Molteplici porte di esecuzione parallele ($IPC > 1$) & Interleaving di accumuli indipendenti nel calcolo di BM25 \\
\textbf{Data Hazards} & Dipendenze seriali strette (RAW) & Spezzamento delle catene di accumulo tramite registri multipli \\
\textbf{Branch Prediction} & Svuotamento della pipeline se il salto fallisce & Codice branchless nelle intersezioni AND e nel pruning WAND \\
\textbf{Vettorizzazione SIMD} & Calcolo simultaneo su registri fino a 512 bit & Decodifica vettoriale di liste compresse (SIMD-BP128, PForDelta) \\
\bottomrule
\end{tabular}%
}
\caption{Quadro sinottico: dai vincoli micro-architetturali della CPU alle scelte ingegneristiche in Information Retrieval.}
\label{tab:sinottico-hardware-ir}
\end{table}
